\section{Xây dựng tập dữ liệu}
\label{sec:dataset}

\subsection{Nguồn dữ liệu}

Tập dữ liệu được xây dựng từ hai nghiên cứu gần đây:

\begin{itemize}
    \item \textbf{DeePhage dataset} \cite{wu2021deephage}: Tập dữ liệu gốc với các phage đã được phân loại lối sống
    \item \textbf{DeepPL annotations} \cite{zhang2024deeppl}: Các chuỗi mới được gán nhãn bổ sung
\end{itemize}

Tập dữ liệu kết hợp bao gồm \textbf{2,241 bộ gen phage hoàn chỉnh}:
\begin{itemize}
    \item 707 phage temperate (nhãn 0)
    \item 1,534 phage virulent (nhãn 1)
    \item Tỷ lệ virulent:temperate $\approx$ 2:1 (class imbalance)
\end{itemize}

\subsection{Tạo contigs bằng sliding window}

Để mô phỏng kịch bản metagenomic thực tế, contigs được tạo từ complete genomes bằng phương pháp sliding window:

\subsubsection{Thuật toán}

\begin{algorithm}[h]
\caption{Contig Generation với Sliding Window}
\label{alg:contig_generation}
\begin{algorithmic}[1]
\State \textbf{Input:} Complete genome $G$ với độ dài $L$, nhóm độ dài $\text{group} \in \{A, B, C, D\}$
\State \textbf{Output:} Tập contigs $\mathcal{C}$
\State $\mathcal{C} \leftarrow \emptyset$
\State $\text{pos} \leftarrow 0$
\While{$\text{pos} < L$}
    \State $\text{length} \sim \mathcal{N}(\mu_{\text{group}}, \sigma_{\text{group}}^2)$ \Comment{Sample từ phân phối chuẩn}
    \State $\text{length} \leftarrow \text{clip}(\text{length}, \text{min}_{\text{group}}, \text{max}_{\text{group}})$
    \State $\text{contig} \leftarrow G[\text{pos}:\text{pos}+\text{length}]$
    \State $\mathcal{C} \leftarrow \mathcal{C} \cup \{\text{contig}\}$
    \State $\text{overlap} \leftarrow \text{length} \times \text{overlap\_ratio}_{\text{group}}$
    \State $\text{pos} \leftarrow \text{pos} + \text{length} - \text{overlap}$
\EndWhile
\State \Return $\mathcal{C}$
\end{algorithmic}
\end{algorithm}

\subsubsection{Tham số cho từng nhóm}

\begin{table}[h]
\centering
\caption{Tham số tạo contig cho các nhóm độ dài}
\label{tab:contig_params}
\begin{tabular}{lcccc}
\toprule
\textbf{Nhóm} & \textbf{Phạm vi (bp)} & \textbf{$\mu$ (bp)} & \textbf{$\sigma$ (bp)} & \textbf{Overlap ratio} \\
\midrule
A & 100-400 & 250 & 75 & 10\% \\
B & 400-800 & 600 & 100 & 20\% \\
C & 800-1200 & 1000 & 100 & 30\% \\
D & 1200-1800 & 1500 & 150 & 40\% \\
\bottomrule
\end{tabular}
\end{table}

\textbf{Rationale cho overlap ratios:}
\begin{itemize}
    \item Nhóm ngắn (A): Overlap thấp để tránh tạo quá nhiều contigs tương tự
    \item Nhóm dài (D): Overlap cao để tăng số lượng samples và đa dạng hóa dữ liệu
\end{itemize}

\subsection{Data augmentation}

\subsubsection{Reverse complement}

Mỗi contig được augment bằng cách tạo reverse complement sequence:

\begin{equation}
\text{RevComp}(\text{ATCG}) = \text{CGAT}
\end{equation}

với quy tắc: $A \leftrightarrow T$, $C \leftrightarrow G$.

\textbf{Lý do:} DNA có cấu trúc sợi đôi, cả hai hướng đều mang thông tin sinh học tương đương. Augmentation này tăng gấp đôi kích thước dataset và cải thiện robustness.

\subsection{Xử lý class imbalance}

\subsubsection{Random undersampling}

Để giải quyết tỷ lệ virulent:temperate = 2:1, áp dụng random undersampling trên training set:

\begin{enumerate}
    \item Giữ nguyên tất cả samples của lớp thiểu số (temperate)
    \item Ngẫu nhiên chọn số lượng samples tương đương từ lớp đa số (virulent)
    \item Kết quả: Training set cân bằng 1:1
\end{enumerate}

\textbf{Trade-off:}
\begin{itemize}
    \item \textit{Ưu điểm}: Đơn giản, hiệu quả, tránh model bias
    \item \textit{Nhược điểm}: Mất một phần thông tin từ lớp đa số
    \item \textit{Lý do chọn}: Giới hạn tài nguyên tính toán, undersampling nhanh hơn các phương pháp phức tạp như SMOTE
\end{itemize}

\subsection{Chia tập dữ liệu}

\subsubsection{Stratified 5-fold cross-validation}

Sử dụng stratified 5-fold CV để đảm bảo:
\begin{itemize}
    \item Tỷ lệ các lớp được duy trì trong mỗi fold
    \item Đánh giá robust trên nhiều splits khác nhau
    \item Giảm variance trong kết quả đánh giá
\end{itemize}

\subsubsection{Train-validation split}

Mỗi fold được chia:
\begin{itemize}
    \item \textbf{Training set}: 80\% (sau undersampling)
    \item \textbf{Validation set}: 20\%
    \item \textbf{Test set}: Fold còn lại (20\% của toàn bộ data)
\end{itemize}

\subsection{Thống kê tập dữ liệu}

\begin{table}[h]
\centering
\caption{Thống kê số lượng contigs sau augmentation (trước undersampling)}
\label{tab:dataset_stats}
\begin{tabular}{lrrrr}
\toprule
\textbf{Nhóm} & \textbf{Temperate} & \textbf{Virulent} & \textbf{Total} & \textbf{Ratio (V:T)} \\
\midrule
A (100-400 bp) & 45,234 & 98,567 & 143,801 & 2.18:1 \\
B (400-800 bp) & 28,156 & 61,423 & 89,579 & 2.18:1 \\
C (800-1200 bp) & 18,945 & 41,289 & 60,234 & 2.18:1 \\
D (1200-1800 bp) & 12,678 & 27,634 & 40,312 & 2.18:1 \\
\bottomrule
\end{tabular}
\end{table}

\textbf{Quan sát:}
\begin{itemize}
    \item Số lượng contigs giảm khi độ dài tăng (do sliding window với overlap)
    \item Tỷ lệ V:T nhất quán ~2.18:1 trên tất cả các nhóm
    \item Sau undersampling, mỗi nhóm có số lượng samples cân bằng 1:1
\end{itemize}
